% Drawing 152/01-15 -- folder 152, batch 01, archive page 1 (PDF page 2).
% Transcription: au-dessous, au crayon : un cercle pâle portant quatre sommets marqués ; de trois d'entre eux part vers l'extérieur un faisceau de trois ou quatre traits ; l'arc entre les sommets de gauche est repassé et fléché.
% Kind: circle with four marked vertices and three bundles of rays. Method 2 (schematic redraw from the facsimile). Status: partial.
% Notes: The faint circle through the four marked vertices (upper left, left, bottom, upper right); the arcs from the upper left vertex to the left one and from the left one to the bottom one are gone over, with an arrowhead pointing down on the first and a small tick on the second; two small points beside the bottom vertex; small arrowheads on the circle on each side of the upper right vertex. Rays, all drawn from their vertex: three from the upper left one, four from the upper right one (one heavy, one curved), three from the bottom one (one heavy). Overlaid on the 400 dpi crop. Still differs: a wavy scribble runs along the heavy ray of the bottom vertex and is not drawn; the curved ray of the upper right vertex leads to a short underlined pencil word that could not be read, and a second illegible pencil word stands at the right of the circle: both left out. Redrawn 2026-11-12 (issue 128): to check again.
% Coordinates: 1 unit = 1 mm = 1/100 of the crop width (bbox in page fractions: [0.6, 0.335, 0.86, 0.5]).
% Redrawn by Claude (claude-fable-5-2), issue #128, 2026-11-12; not checked by a human; not compiled (no TeX in the session).
% Rights: reproduces part of the work -- behind RENDER_TRANSCRIPTS, never published by itself.
\documentclass[tikz,border=3pt]{standalone}
\usepackage{amsmath,amssymb}
\usetikzlibrary{patterns,arrows.meta}
\tikzset{pencil/.style={black!45}, struck/.style={black!35}, mpt/.style={circle,draw,inner sep=1.3pt}, pt/.style={circle,fill,inner sep=0.9pt}, opt/.style={circle,draw,fill=white,inner sep=1.1pt}}
\begin{document}
\begin{tikzpicture}[x=1mm,y=1mm,line cap=round,line join=round,font=\small]
\draw[pencil] plot[smooth, tension=0.55] coordinates {(24,-17.4) (32,-10.5) (41,-9) (51,-11.6) (58.6,-19.6)};
\draw[pencil] plot[smooth, tension=0.55] coordinates {(58.6,-19.6) (61.6,-30) (58,-42) (49,-50.6) (41,-52) (35,-50.6)};
\draw[-{Stealth[length=1.8mm]}] plot[smooth, tension=0.55] coordinates {(24,-17.4) (19,-26) (17.8,-32)};
\draw plot[smooth, tension=0.55] coordinates {(17.8,-32) (17.6,-36) (18.4,-39.4)};
\draw plot[smooth, tension=0.55] coordinates {(18.4,-39.4) (21.4,-44.6) (27,-48) (35,-50.6)};
\draw[thick] (23.2,-45.2) -- (24.2,-48);
\draw (54.6,-13.2) -- (53.4,-15.6);
\draw (59.6,-26) -- (60.6,-24.4) -- (62.2,-25.6);
\node[pt] at (24,-17.4) {};
\node[pt] at (18.4,-39.4) {};
\node[pt] at (35,-50.6) {};
\node[pt] at (58.6,-19.6) {};
\node[pt] at (32.6,-50.3) {};
\node[pt] at (38.2,-51.6) {};
\draw (24,-17.4) -- (15,-3);
\draw (24,-17.4) -- (2,-14.2);
\draw (24,-17.4) -- (5,-22.6);
\draw[thick] (58.6,-19.6) -- (64,-1.6);
\draw plot[smooth, tension=0.55] coordinates {(58.6,-19.6) (64,-10) (70,-5) (84,-2)};
\draw (58.6,-19.6) -- (77,-11.6);
\draw (58.6,-19.6) -- (84,-23);
\draw (35,-50.6) -- (17.6,-71);
\draw (35,-50.6) -- (26,-78);
\draw[thick] (35,-50.6) -- (41.4,-81);
\end{tikzpicture}
\end{document}
